\subsection{3.2 Kiến trúc}

\begin{frame}{Kiến trúc: sáu tầng, thành phần, công nghệ, số dòng}
\muc{3.2 · Kiến trúc}
\HinhChu{slide/3-2_kien_truc_cac_lop.png}{%
\small\begin{itemize}
  \item Giám sát etl\_log 6 bảng; điều phối bằng một lệnh Python
\end{itemize}}
\KetLuan{PostgreSQL, Python, Power BI, GitHub; số bảng và số dòng từng lớp}

\note{[S136] Kiến trúc theo lớp (không vẽ luồng); luồng ở data flow và sơ đồ job. Tiêu đề đầy đủ trước khi rút gọn: Kiến trúc theo lớp: sáu tầng từ nguồn đến BI, mỗi tầng có thành phần, công nghệ, số dòng Tiêu đề trước đây: Kiến trúc sáu tầng từ nguồn đến BI, mỗi tầng có thành phần, công nghệ, số dòng}
\end{frame}

\begin{frame}{Kiến trúc: vai trò từng lớp}
\muc{3.2 · Kiến trúc}
\begin{BangTC}{lXX}
\textbf{Lớp} & \textbf{Việc làm} & \textbf{Số liệu và kiểm tra} \\ \midrule
Nguồn & \texttt{oltp} (9 bảng) và \texttt{crm} (2 bảng) & 11 bảng, 1.559.764 dòng \\
Staging & chép dữ liệu từ nguồn vào \texttt{stg\_*} & staging = ODS + dòng bị loại \\
ODS & làm sạch, gắn cờ \texttt{dq\_flags} & 18 bảng, 678.183 dòng \\
Kho (\texttt{dw}) & 10 dimension và 6 fact & 136 phép kiểm thử \\
Data mart & 4 mart cho 4 dashboard & 17 view, khớp fact \\
Power BI & 4 dashboard & đối chiếu giá trị kỳ vọng từ SQL \\
\end{BangTC}
\KetLuan{Mỗi lớp có một việc và một cách kiểm tra riêng}
\note{[P12] Trang giải thích từng lớp của kiến trúc (checklist: phân tích chi tiết kiến trúc). Số liệu lấy từ các slide ETL và dimension.}
\end{frame}

\begin{frame}{Kiến trúc: công nghệ sử dụng}
\muc{3.2 · Kiến trúc}
\begin{itemize}
  \item Không Docker, không PostGIS; điều phối bằng một lệnh
\end{itemize}
\khungchen[0.3\textheight]{CHÈN HÌNH / BẢNG: logo công nghệ (chèn)}

\note{[S51] Chạy được trên laptop sinh viên 8 GB. Tiêu đề trước đây: Công nghệ: PostgreSQL 18, Python 3.11, SQL procedure, Power BI, Git}
\end{frame}

\subsection{3.3 ETL}

\begin{frame}{Tóm tắt ETL: một lệnh chạy các bước và ghi nhật ký}
\muc{3.3 · ETL}
\HinhLon{slide/3-4-1_pipeline.png}
\KetLuan{extract $\rightarrow$ ODS $\rightarrow$ dim $\rightarrow$ fact $\rightarrow$ mart $\rightarrow$ tests $\rightarrow$ snapshot $\rightarrow$ export}

\note{[S67] Mỗi bước ghi \texttt{job\_\allowbreak{}run}; lỗi thì dừng và đánh dấu failed. Tiêu đề trước đây: Một lệnh chạy sáu bước ETL và ghi lại từng bước}
\end{frame}

\begin{frame}{Tóm tắt ETL: ODS làm sạch 18 bảng, 678.183 dòng}
\muc{3.3 · ETL}
\only<1>{%
\begin{BangTC}{XrX}
\textbf{Lỗi} & \textbf{Số dòng} & \textbf{Xử lý} \\ \midrule
Bàn giao vận chuyển trước khi duyệt thanh toán & 1.359 & Gắn cờ dq\_flags, giữ đơn; khoảng thời gian âm đặt NULL \\
Bàn giao vận chuyển trước ngày mua & 166 & Gắn cờ dq\_flags, giữ đơn; khoảng thời gian âm đặt NULL \\
Giao khách trước khi bàn giao vận chuyển & 23 & Gắn cờ dq\_flags, giữ đơn; khoảng thời gian âm đặt NULL \\
Đơn đã giao nhưng không có ngày duyệt & 14 & Gắn cờ dq\_flags, giữ đơn; khoảng thời gian âm đặt NULL \\
\end{BangTC}
\KetLuan{Zip thành 5 chữ số, thành phố bỏ dấu, \texttt{NULL} thành \texttt{unknown}, khử trùng review}}
\only<2>{%
\begin{BangTC}{XrX}
\textbf{Lỗi} & \textbf{Số dòng} & \textbf{Xử lý} \\ \midrule
Trạng thái delivered nhưng không có ngày giao & 8 & Gắn cờ dq\_flags, giữ đơn; khoảng thời gian âm đặt NULL \\
Có ngày giao nhưng trạng thái chưa là delivered & 6 & Gắn cờ dq\_flags, giữ đơn; khoảng thời gian âm đặt NULL \\
Đã giao nhưng không có ngày bàn giao & 2 & Gắn cờ dq\_flags, giữ đơn; khoảng thời gian âm đặt NULL \\
\end{BangTC}
\KetLuan{Zip thành 5 chữ số, thành phố bỏ dấu, \texttt{NULL} thành \texttt{unknown}, khử trùng review}}
\only<3>{%
\begin{BangTC}{XrX}
\textbf{Lỗi} & \textbf{Số dòng} & \textbf{Xử lý} \\ \midrule
Giá trị thanh toán = 0 & 9 & Gắn cờ, chuẩn hoá (số kỳ 0 → 1) \\
Hình thức thanh toán not\_defined & 3 & Gắn cờ, chuẩn hoá (số kỳ 0 → 1) \\
Số kỳ trả góp = 0 (đặt thành 1) & 2 & Gắn cờ, chuẩn hoá (số kỳ 0 → 1) \\
\end{BangTC}
\KetLuan{Zip thành 5 chữ số, thành phố bỏ dấu, \texttt{NULL} thành \texttt{unknown}, khử trùng review}}
\only<4>{%
\begin{BangTC}{XrX}
\textbf{Lỗi} & \textbf{Số dòng} & \textbf{Xử lý} \\ \midrule
Doanh thu tháng khai báo = 0 & 797 & Gắn cờ, giữ nguyên giá trị \\
Chốt deal trước ngày liên hệ đầu & 1 & Gắn cờ, giữ nguyên giá trị \\
coordinates outside Brazil bounding box & 42 & Loại khỏi ODS, ghi etl\_log.rejected\_rows \\
\end{BangTC}
\KetLuan{Zip thành 5 chữ số, thành phố bỏ dấu, \texttt{NULL} thành \texttt{unknown}, khử trùng review}}
\note{[S52] 1.411 đơn (1,4\%) có ít nhất một cờ chất lượng. Bấm để chuyển giữa 4 phần của bảng. Tiêu đề trước đây: ODS là lớp làm sạch trung gian: 18 bảng, 678.183 dòng, mọi lỗi có cờ hoặc lý do}
\end{frame}

\begin{frame}{Tóm tắt ETL: ODS xử lý lỗi theo ba cách}
\muc{3.3 · ETL}
\only<1>{%
\begin{BangTC}{XrX}
\textbf{Lỗi} & \textbf{Số dòng} & \textbf{Xử lý} \\ \midrule
Bàn giao vận chuyển trước khi duyệt thanh toán & 1.359 & Gắn cờ dq\_flags, giữ đơn; khoảng thời gian âm đặt NULL \\
Bàn giao vận chuyển trước ngày mua & 166 & Gắn cờ dq\_flags, giữ đơn; khoảng thời gian âm đặt NULL \\
Giao khách trước khi bàn giao vận chuyển & 23 & Gắn cờ dq\_flags, giữ đơn; khoảng thời gian âm đặt NULL \\
Đơn đã giao nhưng không có ngày duyệt & 14 & Gắn cờ dq\_flags, giữ đơn; khoảng thời gian âm đặt NULL \\
\end{BangTC}
\KetLuan{1.359 đơn bàn giao trước khi duyệt được gắn cờ, không xóa}}
\only<2>{%
\begin{BangTC}{XrX}
\textbf{Lỗi} & \textbf{Số dòng} & \textbf{Xử lý} \\ \midrule
Trạng thái delivered nhưng không có ngày giao & 8 & Gắn cờ dq\_flags, giữ đơn; khoảng thời gian âm đặt NULL \\
Có ngày giao nhưng trạng thái chưa là delivered & 6 & Gắn cờ dq\_flags, giữ đơn; khoảng thời gian âm đặt NULL \\
Đã giao nhưng không có ngày bàn giao & 2 & Gắn cờ dq\_flags, giữ đơn; khoảng thời gian âm đặt NULL \\
\end{BangTC}
\KetLuan{1.359 đơn bàn giao trước khi duyệt được gắn cờ, không xóa}}
\only<3>{%
\begin{BangTC}{XrX}
\textbf{Lỗi} & \textbf{Số dòng} & \textbf{Xử lý} \\ \midrule
Giá trị thanh toán = 0 & 9 & Gắn cờ, chuẩn hoá (số kỳ 0 → 1) \\
Hình thức thanh toán not\_defined & 3 & Gắn cờ, chuẩn hoá (số kỳ 0 → 1) \\
Số kỳ trả góp = 0 (đặt thành 1) & 2 & Gắn cờ, chuẩn hoá (số kỳ 0 → 1) \\
\end{BangTC}
\KetLuan{1.359 đơn bàn giao trước khi duyệt được gắn cờ, không xóa}}
\only<4>{%
\begin{BangTC}{XrX}
\textbf{Lỗi} & \textbf{Số dòng} & \textbf{Xử lý} \\ \midrule
Doanh thu tháng khai báo = 0 & 797 & Gắn cờ, giữ nguyên giá trị \\
Chốt deal trước ngày liên hệ đầu & 1 & Gắn cờ, giữ nguyên giá trị \\
coordinates outside Brazil bounding box & 42 & Loại khỏi ODS, ghi etl\_log.rejected\_rows \\
\end{BangTC}
\KetLuan{1.359 đơn bàn giao trước khi duyệt được gắn cờ, không xóa}}
\note{[S69] Khoảng thời gian âm đặt NULL ở fact và được đếm lại. Bấm để chuyển giữa 4 phần của bảng. Tiêu đề trước đây: ODS xử lý lỗi theo ba cách: chuẩn hóa, gắn cờ hoặc loại có lý do}
\end{frame}

\begin{frame}{ETL jobs: sáu giai đoạn, 142 test đạt}
\muc{3.3 · ETL}
\HinhChu{slide/3-3_etl_jobs.png}{%
\small\begin{itemize}
  \item Số dòng lấy từ etl\_log.job\_run (lô đầu và tổng các lô)
\end{itemize}}
\KetLuan{Extract, seed và ODS, dimension và fact, mart, kiểm thử, xuất CSV}

\note{[S137] Sơ đồ các job của run_pipeline sinh từ etl_log. Tiêu đề đầy đủ trước khi rút gọn: Sáu giai đoạn ETL: lô đầu nạp 1.273.620 dòng vào staging, 42 dòng loại, 142 test đạt Tiêu đề trước đây: Sáu giai đoạn ETL: lô đầu 1.273.620 dòng vào staging, 42 dòng loại, 142 test đạt}
\end{frame}

\begin{frame}{ETL jobs: làm giàu khoảng cách haversine}
\muc{3.3 · ETL}
\begin{itemize}
  \item 555 dòng (0,49\%) thiếu toạ độ vẫn được giữ lại với giá trị trống
\end{itemize}
\khungchen[0.3\textheight]{CHÈN HÌNH / BẢNG: đoạn SQL \texttt{dw.haversine\_\allowbreak{}km} (Phụ lục B)}

\note{[S107] São Paulo – Rio tính ra 360,9 km, gần với thực tế. Tiêu đề đầy đủ trước khi rút gọn: Mở rộng địa lý: hàm haversine SQL thuần tính khoảng cách seller $\rightarrow$ khách cho 112.650 dòng hàng Tiêu đề trước đây: Haversine SQL thuần tính khoảng cách seller $\rightarrow$ khách cho 112.650 dòng hàng}
\end{frame}

\begin{frame}{Before – after: bảng đơn hàng}
\muc{3.3 · ETL}
\begin{BangTC}{llllX}
\textbf{order\_id} & \textbf{approved\_at} & \textbf{carrier} & \textbf{customer} & \textbf{dq\_flags (sau)} \\ \midrule
0184d4dd\ldots{} & 2017-09-13 22:17:15 & 2017-09-04 14:05:50 & 2017-09-09 15:12:44 & \texttt{carrier\_\allowbreak{}before\_\allowbreak{}approved} \\
0922ee16\ldots{} & 2017-07-11 10:55:15 & 2017-07-14 12:41:34 & 2017-07-12 20:49:42 & \texttt{delivered\_\allowbreak{}before\_\allowbreak{}carrier} \\
12a95a3c\ldots{} & NULL & 2017-02-22 11:23:11 & 2017-03-02 11:09:19 & \texttt{delivered\_\allowbreak{}without\_\allowbreak{}approval} \\
2aa91108\ldots{} & 2017-09-29 09:07:16 & NULL & 2017-11-20 19:44:47 & \texttt{delivered\_\allowbreak{}without\_\allowbreak{}carrier\_\allowbreak{}date} \\
\end{BangTC}
\KetLuan{Ví dụ trước/sau}

\note{[S70] Cùng 99.441 đơn nhưng thêm \texttt{dq\_\allowbreak{}flags}. Tiêu đề trước đây: Trước và sau ETL: bảng đơn hàng có thêm cờ lỗi và kiểu dữ liệu chuẩn}
\end{frame}

\begin{frame}{Before – after: 1.000.163 dòng toạ độ thành 19.178 zip}
\muc{3.3 · ETL}
\TruocSau{Trước (staging)}{zip 01006: 9 điểm\\lat từ $-23{,}55070$ đến $-23{,}55026$}{Sau (ods)}{1 dòng: lat $-23{,}5503$\\lng $-46{,}6362$, sao paulo, SP}
\KetLuan{Trung vị toạ độ, mode thành phố}

\note{[S72] 42 điểm ngoài Brazil được ghi vào bảng loại. Tiêu đề trước đây: Trước và sau ETL: 1.000.163 dòng toạ độ gom thành 19.178 mã bưu chính}
\end{frame}

\begin{frame}{OLTP $\rightarrow$ OLAP: bảng nguồn và bảng đích}
\muc{3.3 · ETL}
\only<1>{%
\begin{BangTC}{XlXX}
\textbf{Bảng nguồn} & \textbf{Cột dùng} & \textbf{Bảng đích} & \textbf{Phép biến đổi tiêu biểu} \\ \midrule
oltp.customers & 5/5 & dim\_customer, dim\_geography và 4 fact & zip $\rightarrow$ char(5), lpad bằng số 0 \\
oltp.geolocation & 5/6 & dim\_geography & chuẩn hoá char(5) rồi GROUP BY zip \\
oltp.order\_items & 7/7 & fact\_order\_fulfillment, fact\_order\_item & giữ nguyên giá trị \\
\end{BangTC}
\KetLuan{65 trong 71 cột nguồn được dùng sang OLAP}}
\only<2>{%
\begin{BangTC}{XlXX}
\textbf{Bảng nguồn} & \textbf{Cột dùng} & \textbf{Bảng đích} & \textbf{Phép biến đổi tiêu biểu} \\ \midrule
oltp.order\_payments & 5/5 & fact\_order\_fulfillment, fact\_payment & payment\_type: lower(trim()) \\
oltp.order\_reviews & 7/7 & fact\_review & khử trùng cặp (review\_id, order\_id), giữ bản mới nhất \\
oltp.orders & 8/8 & dim\_customer và 4 fact & order\_status: lower(trim()) \\
\end{BangTC}
\KetLuan{65 trong 71 cột nguồn được dùng sang OLAP}}
\only<3>{%
\begin{BangTC}{XlXX}
\textbf{Bảng nguồn} & \textbf{Cột dùng} & \textbf{Bảng đích} & \textbf{Phép biến đổi tiêu biểu} \\ \midrule
oltp.product\_category\_name\_translation & 2/2 & dim\_product & giữ nguyên giá trị \\
oltp.products & 6/9 & dim\_product & category rỗng $\rightarrow$ 'unknown' \\
oltp.sellers & 4/4 & dim\_geography, dim\_seller & zip $\rightarrow$ char(5) \\
\end{BangTC}
\KetLuan{65 trong 71 cột nguồn được dùng sang OLAP}}
\only<4>{%
\begin{BangTC}{XlXX}
\textbf{Bảng nguồn} & \textbf{Cột dùng} & \textbf{Bảng đích} & \textbf{Phép biến đổi tiêu biểu} \\ \midrule
crm.closed\_deals & 12/14 & dim\_lead\_profile, dim\_sales\_rep, dim\_seller, fact\_lead\_funnel & business\_segment: trim, NULL $\rightarrow$ 'unknown' \\
crm.mql & 4/4 & dim\_seller, fact\_lead\_funnel & origin: lower(trim()), NULL $\rightarrow$ 'unknown' \\
\end{BangTC}
\KetLuan{65 trong 71 cột nguồn được dùng sang OLAP}}
\note{[P12] Tổng hợp mapping OLTP sang OLAP theo bảng nguồn (báo cáo mục 3.3, bảng tab:mapping-tong-hop; mapping đầy đủ từng cột ở Phụ lục C). Cột dùng là số cột dùng sang OLAP trên tổng số cột của bảng nguồn. Bấm để chuyển giữa hai phần.}
\end{frame}

\begin{frame}{Đa dạng công cụ ETL: SQL, Python, Power Query}
\muc{3.3 · ETL}
% cho-phep: R10 ba công cụ đại diện, chi tiết ở ghi chú
\begin{itemize}
  \item Chữ thường, bỏ dấu, zip đủ 5 ký tự; 73.440 dòng đổi
  \item Python chênh 2 dòng do ký tự ordinal; mã M thêm bảng thay thế để khớp SQL
  \item Mã M và hướng dẫn chụp before/after: powerbi/power-query-vi-du.md
\end{itemize}

\note{[S138] Nhóm cần chạy mã M trong Power BI và chụp ảnh before/after. Tiêu đề đầy đủ trước khi rút gọn: SQL, Python và Power Query cho cùng kết quả làm sạch: 8.011 thành 5.968 thành phố Tiêu đề trước đây: SQL, Python và Power Query cùng kết quả làm sạch: 8.011 thành 5.968 thành phố}
\end{frame}

\begin{frame}{Đối soát số dòng: nguồn = staging = ODS + loại}
\muc{3.3 · ETL}
\begin{itemize}
  \item $\Sigma$ price 13.591.643,70 BRL = file CSV
\end{itemize}
\khungchen[0.3\textheight]{CHÈN HÌNH / BẢNG: bảng đối soát (handoff P4)}

\note{[S74] Chênh lệch nếu có phải giải thích được bằng \texttt{rejected\_\allowbreak{}rows}. Tiêu đề trước đây: Đối soát số dòng: nguồn = staging = ODS + loại, và tổng tiền khớp từng đồng}
\end{frame}

\begin{frame}{Nạp tăng dần: 11 đợt, 286.144 bản ghi mới}
\muc{3.3 · ETL}
\HinhLon{slide/3-4-1_pipeline.png}
\KetLuan{\texttt{simulate\_\allowbreak{}increment} rồi \texttt{run\_\allowbreak{}pipeline --mode incremental}}

\note{[S77] Đơn vào với trạng thái cuối; đó là giới hạn của mô phỏng. Tiêu đề trước đây: Nạp tăng dần mô phỏng: 11 đợt hằng tháng, tổng 286.144 bản ghi mới}
\end{frame}

\begin{frame}{Nạp tăng dần: incremental trùng nạp full, 16/16 bảng}
\muc{3.3 · ETL}
\begin{itemize}
  \item So sánh bằng checksum trên khoá tự nhiên
\end{itemize}
\khungchen[0.3\textheight]{CHÈN HÌNH / BẢNG: bảng so sánh (\texttt{handoff/2026-11-05-P5.md})}

\note{[S80] Khoá thay thế khác nhau nên so sánh qua khoá tự nhiên. Tiêu đề đầy đủ trước khi rút gọn: Kết quả incremental trùng nạp full: 16/16 bảng khớp số dòng, tổng tiền và nội dung Tiêu đề trước đây: Incremental trùng nạp full: 16/16 bảng khớp số dòng, tổng tiền và nội dung}
\end{frame}

\begin{frame}{Kiểm thử ETL: 136 phép đạt 100\%, cộng 6 phép đối chiếu}
\muc{3.3 · ETL}
\begin{BangTC}{Xr}
\textbf{Nhóm kiểm thử} & \textbf{Số phép} \\ \midrule
ODS (số dòng, loại, cleaning) & 25 \\
Dimension (Unknown, SCD2, phủ khoá) & 42 \\
Fact (đối soát, grain, khoá ngoại, logic) & 49 \\
Data mart (khớp fact) & 17 \\
Thử lại dòng mồ côi (kịch bản) & 3 \\
\textbf{Tổng} & \textbf{136 (đạt 100\%)} \\
Đối chiếu khám phá dữ liệu (chỉ báo cáo) & 6/6 khớp \\
\end{BangTC}
\KetLuan{Số dòng, tổng tiền, grain, khoá ngoại, SCD2, mart khớp fact}
\note{[S82] Chạy sau mỗi lần pipeline, lưu \texttt{etl\_\allowbreak{}log.dq\_\allowbreak{}check}. Tiêu đề trước đây: 136 phép kiểm thử đạt 100\%, cộng 6 phép đối chiếu với khám phá dữ liệu}
\end{frame}

\begin{frame}{Dựng lại từ đầu: 443 giây, kết quả giống hệt}
\muc{3.3 · ETL}
\only<1>{%
\begin{BangTC}{XrX}
\textbf{Bước} & \textbf{Giây} & \textbf{Kết quả} \\ \midrule
drop olist\_dw + olist\_full & 0.6 & dropped \\
setup\_db (schemas, DDL, procedures, views) & 1.8 & Setup done. \\
load\_source --initial (olist\_dw, to 2017-12-31) & 15.5 & Rejected rows logged in etl\_log.rejected\_rows: 0 (batch 1) \\
load\_source --full (olist\_full) & 24.0 & Rejected rows logged in etl\_log.rejected\_rows: 0 (batch 1) \\
\end{BangTC}
\KetLuan{Xóa hai database, chạy một lệnh, 16/16 bảng khớp}}
\only<2>{%
\begin{BangTC}{XrX}
\textbf{Bước} & \textbf{Giây} & \textbf{Kết quả} \\ \midrule
run\_pipeline --mode full (initial state) & 55.8 & tests: 136/136 đạt, 0 lỗi; đối chiếu EDA 6/6 khớp; pipeline 54.7s \\
replay\_increments (11 x simulate + incremental) & 287.1 & 11 đợt hằng tháng (01/2018 → 11/2018); sau đợt cuối: 8,000 MQL, 842 deal, 379 seller có đơn bán được, 99,441 đơn, GMV 13,494,400.74 \\
run\_pipeline --mode incremental (tests + export\_powerbi) & 30.5 & tests: 136/136 đạt, 0 lỗi; đối chiếu EDA 6/6 khớp; pipeline 29.3s \\
\end{BangTC}
\KetLuan{Xóa hai database, chạy một lệnh, 16/16 bảng khớp}}
\note{[S110] Khả năng tái lập là bằng chứng độ tin cậy của pipeline. Bấm để chuyển giữa 2 phần của bảng. Tiêu đề trước đây: Dựng lại toàn bộ dự án từ đầu mất 443 giây và cho kết quả giống hệt}
\end{frame}

\subsection{3.4 Dimension}

\begin{frame}{Số dòng mỗi dimension: từ 6 đến 96.339}
\muc{3.4 · Dimension}
\only<1>{%
\begin{BangTC}{lXrc}
\textbf{Chiều} & \textbf{Grain} & \textbf{Số dòng} & \textbf{SCD} \\ \midrule
dim\_date & 1 ngày (2016-01-01 → 2020-12-31) & 1.828 & -- \\
dim\_geography & 1 mã bưu chính (zip) & 19.178 & 1 \\
dim\_customer & 1 phiên bản địa chỉ của 1 khách thật & 96.339 & 2 \\
dim\_seller & 1 seller (chiều dùng chung) & 3.558 & 1 \\
dim\_product & 1 sản phẩm & 32.952 & 1 \\
\end{BangTC}
\KetLuan{dim\_date 1.828; geography 19.178; customer 96.339; seller 3.558; product 32.952; ...}}
\only<2>{%
\begin{BangTC}{lXrc}
\textbf{Chiều} & \textbf{Grain} & \textbf{Số dòng} & \textbf{SCD} \\ \midrule
dim\_order\_status & 1 trạng thái đơn & 9 & 1 \\
dim\_payment\_type & 1 hình thức thanh toán & 6 & 1 \\
dim\_lead\_origin & 1 giá trị origin của MQL & 11 & 1 \\
dim\_lead\_profile & 1 tổ hợp thuộc tính deal (junk dimension) & 503 & -- \\
dim\_sales\_rep & 1 nhân viên và vai trò (SDR/SR) & 55 & -- \\
\end{BangTC}
\KetLuan{dim\_date 1.828; geography 19.178; customer 96.339; seller 3.558; product 32.952; ...}}
\note{[S57] Mỗi chiều có đúng một dòng Unknown, được kiểm thử. Bấm để chuyển giữa 2 phần của bảng. Tiêu đề trước đây: Mười bảng chiều với số dòng từ 6 đến 96.339}
\end{frame}

\begin{frame}{Nhóm dimension: sáu nhóm theo vai trò}
\muc{3.4 · Dimension}
\begin{BangTC}{lXr}
\textbf{Nhóm} & \textbf{Dimension} & \textbf{Số dòng} \\ \midrule
Dùng chung & \texttt{dim\_seller} & 3.558 \\
Đóng vai (role-playing) & \texttt{dim\_date}, \texttt{dim\_geography} & 1.828; 19.178 \\
SCD2 & \texttt{dim\_customer} & 96.339 \\
Chiều rác (junk) & \texttt{dim\_lead\_profile} & 503 \\
Danh mục nhỏ & \texttt{dim\_order\_status}, \texttt{dim\_payment\_type}, \texttt{dim\_lead\_origin}, \texttt{dim\_sales\_rep} & 9; 6; 11; 55 \\
Sản phẩm & \texttt{dim\_product} & 32.952 \\
\end{BangTC}
\KetLuan{Mười chiều chia nhóm: dùng chung, đóng vai, SCD2, rác, danh mục nhỏ}
\note{[P12] Phân dimension thành nhóm để dễ theo dõi (checklist). Số dòng lấy từ bảng số dòng mỗi dimension.}
\end{frame}

\begin{frame}{Hệ thống chiều khái niệm: phân cấp năm chiều}
\muc{3.4 · Dimension}
\begin{itemize}
  \item Ngày $\rightarrow$ tháng $\rightarrow$ quý $\rightarrow$ năm; zip $\rightarrow$ thành phố $\rightarrow$ bang $\rightarrow$ vùng
\end{itemize}
\khungchen[0.3\textheight]{CHÈN HÌNH / BẢNG: sơ đồ phân cấp (vẽ)}

\note{[S58] Phân cấp là cơ sở của roll-up và drill-down. Tiêu đề trước đây: Hệ thống chiều khái niệm: phân cấp ngày, địa lý, sản phẩm, seller, khách}
\end{frame}

\begin{frame}{Hệ thống chiều khái niệm: số giá trị mỗi cấp}
\muc{3.4 · Dimension}
\begin{itemize}
  \item 12 số tháng nhưng 60 cặp (năm, tháng); 5.814 tên thành phố nhưng 6.190 cặp (bang, thành phố)
  \item Mỗi danh mục một nhóm ngành, mỗi origin một nhóm kênh (0 vi phạm)
  \item Mười chiều chia nhóm: dùng chung, đóng vai, SCD2, rác, danh mục nhỏ
\end{itemize}

\note{[S139] Bảng phân cấp từng chiều với giá trị ví dụ và số giá trị phân biệt ở báo cáo mục 3.4. Tiêu đề đầy đủ trước khi rút gọn: Phân cấp đúng khi tháng kèm năm, thành phố kèm bang: 303 tên thành phố trùng giữa các bang Tiêu đề trước đây: Phân cấp đúng: tháng kèm năm, thành phố kèm bang; 303 tên thành phố trùng}
\end{frame}

\begin{frame}{Dimension: SCD2 địa chỉ khách, 96.339 dòng}
\muc{3.4 · Dimension}
\only<1>{%
\begin{BangTC}{lllX}
\textbf{customer\_key} & \textbf{khách thật} & \textbf{zip} & \textbf{thành phố} \\ \midrule
15081 & 570eb70f… & 31930 & belo horizonte \\
15082 & 570eb70f… & 31130 & belo horizonte \\
\end{BangTC}
\KetLuan{Đơn cũ vẫn trỏ về địa chỉ lúc mua}}
\only<2>{%
\begin{BangTC}{lXXXX}
\textbf{customer\_key} & \textbf{effective\_from} & \textbf{effective\_to} & \textbf{is\_current} & \textbf{version} \\ \midrule
15081 & 1900-01-01 & 2017-11-08 & không & 1 \\
15082 & 2017-11-09 & 9999-12-31 & có & 2 \\
\end{BangTC}
\KetLuan{Đơn cũ vẫn trỏ về địa chỉ lúc mua}}
\note{[S59] Ví dụ khách \texttt{570eb70f…} đổi zip ngày 09/11/2017; đơn 24/10 thuộc phiên bản 1. Bấm để chuyển giữa 2 phần của bảng. Tiêu đề trước đây: SCD2 giữ lịch sử địa chỉ: 96.096 khách thành 96.339 dòng, 236 khách đổi zip}
\end{frame}

\subsection{3.5 Data model OLAP}

\begin{frame}{Data model logic: bus matrix sáu fact}
\muc{3.5 · Data model OLAP}
\only<1>{%
\begin{BangTC}{Xcccc}
\textbf{Fact} & \textbf{Date} & \textbf{Customer} & \textbf{Seller} & \textbf{Product} \\ \midrule
fact\_order\_item & \ding{51} & \ding{51} & \ding{51} & \ding{51} \\
fact\_order\_fulfillment & \ding{51} & \ding{51} &  &  \\
fact\_payment & \ding{51} & \ding{51} &  &  \\
fact\_review & \ding{51} & \ding{51} &  &  \\
fact\_lead\_funnel & \ding{51} &  & \ding{51} &  \\
fact\_seller\_monthly & \ding{51} &  & \ding{51} &  \\
\end{BangTC}
\KetLuan{Seller là chiều dùng chung của hai quy trình}}
\only<2>{%
\begin{BangTC}{Xcccc}
\textbf{Fact} & \textbf{Geo} & \textbf{Status} & \textbf{Payment} & \textbf{Origin} \\ \midrule
fact\_order\_item & \ding{51} & \ding{51} &  &  \\
fact\_order\_fulfillment & \ding{51} & \ding{51} &  &  \\
fact\_payment &  &  & \ding{51} &  \\
fact\_review &  &  &  &  \\
fact\_lead\_funnel &  &  &  & \ding{51} \\
fact\_seller\_monthly &  &  &  &  \\
\end{BangTC}
\KetLuan{Seller là chiều dùng chung của hai quy trình}}
\only<3>{%
\begin{BangTC}{Xcc}
\textbf{Fact} & \textbf{Profile} & \textbf{Rep} \\ \midrule
fact\_order\_item &  &  \\
fact\_order\_fulfillment &  &  \\
fact\_payment &  &  \\
fact\_review &  &  \\
fact\_lead\_funnel & \ding{51} & \ding{51} \\
fact\_seller\_monthly &  &  \\
\end{BangTC}
\KetLuan{Seller là chiều dùng chung của hai quy trình}}
\note{[S53] Bus matrix cho thấy chiều nào được dùng lại ở đâu. Bấm để chuyển giữa 3 phần của bảng. Tiêu đề trước đây: Bus matrix: sáu fact dùng chung các chiều ngày, seller, khách, địa lý}
\end{frame}

\begin{frame}{Data model logic: constellation bốn mart}
\muc{3.5 · Data model OLAP}
\only<1>{%
\HinhLon{slide/3-3-4_constellation_a.png}}\only<2>{%
\HinhLon{slide/3-3-4_constellation_b.png}}
\KetLuan{Khung = data mart}

\note{[S54] Fact màu cam, chiều màu xanh đúng yêu cầu trình bày. Tiêu đề trước đây: Mô hình constellation: bốn mart, sáu fact (cam), mười chiều (xanh)}
\end{frame}

\begin{frame}{Data model vật lý: sinh từ cơ sở dữ liệu}
\muc{3.5 · Data model OLAP}
\HinhLon{slide/3-5_data_model_vat_ly.png}
\KetLuan{16 bảng, fact và dimension khác màu}
\note{[S141] Mô hình vật lý tổng thể. Tiêu đề trước đây: Data model vật lý sinh từ cơ sở dữ liệu: mọi bảng dw với cột, kiểu, PK và FK}
\end{frame}

\begin{frame}{Data model vật lý: khoá chính, khoá ngoại, index}
\muc{3.5 · Data model OLAP}
\only<1>{%
\HinhLon{slide/3-3-4_erd_dw_fact_order_item.png}}
\only<2>{%
\HinhLon{slide/3-3-4_erd_dw_fact_order_fulfillment.png}}
\only<3>{%
\HinhLon{slide/3-3-4_erd_dw_fact_payment.png}}
\only<4>{%
\HinhLon{slide/3-3-4_erd_dw_fact_review.png}}
\only<5>{%
\HinhLon{slide/3-3-4_erd_dw_fact_lead_funnel.png}}
\only<6>{%
\HinhLon{slide/3-3-4_erd_dw_fact_seller_monthly.png}}
\KetLuan{Mọi khoá ngoại NOT NULL, dòng Unknown -1}

\note{[S56] Tiền \texttt{numeric(12,2)}, khoá \texttt{<tên>\_\allowbreak{}key}, ngày \texttt{yyyymmdd}. Tiêu đề trước đây: Data model vật lý: khoá chính theo grain, khoá ngoại có index, 112.650 dòng hàng}
\end{frame}

\begin{frame}{Data model: ba loại fact của Kimball}
\muc{3.5 · Data model OLAP}
\begin{BangTC}{llXr}
\textbf{Fact} & \textbf{Loại} & \textbf{Grain} & \textbf{Số dòng} \\ \midrule
fact\_order\_item & Transaction & 1 dòng hàng của 1 đơn & 112.650 \\
fact\_order\_fulfillment & Accumulating snapshot & 1 đơn (5 mốc thời gian) & 99.441 \\
fact\_payment & Transaction & 1 lần thanh toán & 103.886 \\
fact\_review & Transaction & 1 review của 1 đơn & 99.224 \\
fact\_lead\_funnel & Accumulating snapshot & 1 MQL & 8.000 \\
fact\_seller\_monthly & Periodic snapshot (dày đặc) & 1 seller × 1 tháng & 36.583 \\
\end{BangTC}
\KetLuan{Transaction $\times$3, accumulating $\times$2, periodic $\times$1}
\note{[S62] Tổng 459.784 dòng fact (99.441 + 112.650 + 103.886 + 99.224 + 8.000 + 36.583). Tiêu đề trước đây: Sáu fact thuộc đủ ba loại của Kimball}
\end{frame}

\begin{frame}{Data mart: mỗi mart một nhóm câu hỏi}
\muc{3.5 · Data model OLAP}
\begin{BangTC}{lXX}
\textbf{Mart} & \textbf{View chính} & \textbf{Câu hỏi} \\ \midrule
\texttt{mart\_\allowbreak{}sales} & \texttt{v\_\allowbreak{}sales\_\allowbreak{}daily}, \texttt{v\_\allowbreak{}sales\_\allowbreak{}state\_\allowbreak{}month}, \texttt{v\_\allowbreak{}category\_\allowbreak{}pareto} & GMV theo ngày, bang, danh mục; Pareto \\
\texttt{mart\_\allowbreak{}fulfillment} & \texttt{v\_\allowbreak{}fulfillment\_\allowbreak{}order}, \texttt{v\_\allowbreak{}fulfillment\_\allowbreak{}state\_\allowbreak{}month}, \texttt{v\_\allowbreak{}shipping\_\allowbreak{}distance} & lead time, trễ, trạng thái đơn, khoảng cách \\
\texttt{mart\_\allowbreak{}cx} & \texttt{v\_\allowbreak{}review\_\allowbreak{}vs\_\allowbreak{}delay}, \texttt{v\_\allowbreak{}payment\_\allowbreak{}mix\_\allowbreak{}month} & review theo mức trễ; cơ cấu thanh toán \\
\texttt{mart\_\allowbreak{}acquisition} & \texttt{v\_\allowbreak{}funnel\_\allowbreak{}by\_\allowbreak{}origin\_\allowbreak{}month} & chuyển đổi MQL theo kênh \\
\end{BangTC}
\KetLuan{mart\_sales, mart\_fulfillment, mart\_cx, mart\_acquisition}
\note{[S142] Data mart. Tiêu đề trước đây: Bốn data mart, mỗi mart gồm fact, dimension và 3--4 câu hỏi phân tích}
\end{frame}

\begin{frame}{OLAP: roll-up và drill-down}
\muc{3.5 · Data model OLAP}
\only<1>{%
\begin{BangTC}{llrr}
\textbf{Năm} & \textbf{Quý} & \textbf{Số đơn} & \textbf{GMV (BRL)} \\ \midrule
2017 & 2017-Q1 & 5.122 & 733.398,94 \\
2017 & 2017-Q2 & 9.222 & 1.286.918,78 \\
2017 & 2017-Q3 & 12.445 & 1.681.949,00 \\
2017 & 2017-Q4 & 17.586 & 2.406.225,55 \\
2017 & (cả năm) & 44.375 & 6.108.492,27 \\
\end{BangTC}
\KetLuan{Quý 4/2017: 2,41 triệu BRL; 24/11/2017: 152.654 BRL}}
\only<2>{%
\begin{BangTC}{llrr}
\textbf{Năm} & \textbf{Quý} & \textbf{Số đơn} & \textbf{GMV (BRL)} \\ \midrule
2018 & 2018-Q1 & 20.979 & 2.764.402,78 \\
2018 & 2018-Q2 & 19.897 & 2.849.730,26 \\
2018 & 2018-Q3 & 12.655 & 1.727.049,37 \\
2018 & (cả năm) & 53.531 & 7.341.182,41 \\
Tổng & (cả năm) & 97.906 & 13.449.674,68 \\
\end{BangTC}
\KetLuan{Quý 4/2017: 2,41 triệu BRL; 24/11/2017: 152.654 BRL}}
\note{[S85] Roll-up dùng \texttt{ROLLUP}, drill-down thêm cấp vào \texttt{GROUP BY}. Tiêu đề đầy đủ trước khi rút gọn: Roll-up và drill-down: GMV 2017 là 6,11 triệu, 2018 là 7,34 triệu; ngày Black Friday cao gấp 5,2 lần Bấm để chuyển giữa 2 phần của bảng. Tiêu đề trước đây: Roll-up: GMV 2017 là 6,11 triệu, 2018 là 7,34 triệu; Black Friday gấp 5,2 lần}
\end{frame}

\begin{frame}{OLAP: slice, dice và pivot}
\muc{3.5 · Data model OLAP}
\begin{BangTC}{Xrrrr}
\textbf{Vùng} & \textbf{2016} & \textbf{2017} & \textbf{2018} & \textbf{Tổng} \\ \midrule
Sudeste & 27.701 & 3.928.343 & 4.866.293 & 8.822.336 \\
Sul & 7.456 & 883.020 & 1.047.636 & 1.938.111 \\
Nordeste & 6.013 & 727.733 & 804.568 & 1.538.314 \\
Centro-Oeste & 2.356 & 401.380 & 458.090 & 861.825 \\
Norte & 1.200 & 168.018 & 164.596 & 333.814 \\
\end{BangTC}
\KetLuan{Dice Sudeste $\times$ Home \& Furniture $\times$ quý 4: SP 199.179 BRL}
\note{[S86] Pivot bằng tổng có điều kiện \texttt{FILTER}. Tiêu đề trước đây: Slice, dice và pivot: vùng Sudeste chiếm 64,3\% GMV năm 2017}
\end{frame}

\begin{frame}{Hiệu năng: star schema nhanh hơn OLTP 1,0–1,6 lần}
\muc{3.5 · Data model OLAP}
\HinhLon{slide/3-5-1_benchmark.png}
\KetLuan{Một seller nhanh gấp 9 lần nhờ index; materialized view dưới 1 ms so với 150–750 ms}

\note{[S88] Không phóng đại: giá trị chính của mô hình chiều là truy vấn đơn giản và số liệu nhất quán. Tiêu đề đầy đủ trước khi rút gọn: Star schema chỉ nhanh hơn OLTP 1,0–1,6 lần; tốc độ lớn đến từ materialized view và index Thanh kết luận đầy đủ trước khi rút gọn: Chọn lọc một seller nhanh gấp 9 lần nhờ index; materialized view dưới 1 ms so với 150–750 ms. Tiêu đề trước đây: Star schema chỉ nhanh hơn OLTP 1,0–1,6 lần; tốc độ nhờ materialized view, index}
\end{frame}
